\chapter{Implementation}

\section{Module Organisation and Separation of Concerns}

The codebase separates mathematical concerns (algorithms) from engineering
concerns (backbone, training, evaluation, UI) using a strict layering policy.
Table~\ref{tab:modules} summarises the core modules.

\begin{table}[htbp]\centering
\caption{Core module responsibilities.}\label{tab:modules}
\begin{tabularx}{\textwidth}{>{\raggedright\arraybackslash}p{.34\textwidth}>{\raggedright\arraybackslash}X}\toprule
Module & Responsibility\\\midrule
\texttt{config/} & Validated configuration dataclass. Algorithm kwargs cannot shadow shared controls\\
\texttt{data/} & CIFAR-10 download, normalisation, dataloader construction\\
\texttt{models/backbone.py} & Sinusoidal time embedding, ResBlocks, SimpleUNet\\
\texttt{algorithms/*.py} & Mathematical objective, JVP target, and sampler logic for each algorithm\\
\texttt{training/trainer.py} & Optimisation loop, AMP, timing, checkpointing\\
\texttt{sampling/sampler.py} & NFE sweep, batched generation, sample-grid saving\\
\texttt{evaluation/} & FID reference cache, FID/IS computation via TorchMetrics\\
\path{experiments/runner.py} & End-to-end orchestration; builds shared objects, runs train/sample/evaluate\\
\path{web/inference_server.py} & HTTP server; auto-discovers run directories; serialises GPU requests with a threading lock\\
\texttt{*.html} & Home, training history, and live-inference browser views\\\bottomrule
\end{tabularx}\end{table}

\section{Algorithm Registry}

All algorithms are registered in \texttt{algorithms/\_\_init\_\_.py} under string
keys (\texttt{fm}, \texttt{fm\_lognorm}, \texttt{mf}, etc.). The CLI
\texttt{train.py} selects an algorithm class from this registry via
\texttt{--algorithm}, and \texttt{evaluate.py} does the same for standalone
evaluation. This pattern ensures no algorithm can be accidentally selected or
skipped.

\section{Flow Matching Implementation}

\texttt{FlowMatchingAlgorithm.training\_step} receives a batch of data images,
samples $t \sim \mathcal{U}(0, 1)$ per image, constructs the linear interpolant
$x_t = (1-t)x_0 + t\varepsilon$, calls the backbone with $(x_t, t)$, and
returns the MSE against the target $\varepsilon - x_0$. \texttt{sample} starts
at pure noise and applies backward Euler, calling the backbone once per step.

\section{Logit-Normal FM Implementation}

\texttt{FlowMatchingLognormAlgorithm} inherits from
\texttt{FlowMatchingAlgorithm} and overrides only the time-sampling line:
\begin{lstlisting}[language=Python]
t = torch.sigmoid(torch.randn(batch_size, device=device))
\end{lstlisting}
The path construction, loss, sampler, and backbone call are identical to FM.
This single-line override is the complete difference between the two algorithms.

\section{Mean Flow Correctness-Critical Path}

The JVP is computed by:
\begin{lstlisting}[language=Python]
tangents = (v_theta, zeros_like(r), ones_like(t))
_, jvp_out = torch.func.jvp(self._forward, (z, r, t), tangents)
target = (v_theta - (t - r) * jvp_out).detach()
\end{lstlisting}
The critical points are:
\begin{itemize}
  \item Tangents are $(v_\theta, 0, 1)$ for $(z, r, t)$, implementing the
        total time derivative $dv/dt$ along the ODE.
  \item The target is \texttt{.detach()}-ed before the MSE loss, preventing
        second-order gradients through the JVP.
  \item AMP autocast is disabled inside \texttt{jvp} to avoid dtype mismatches
        in forward-mode AD.
  \item The $r$-embedding (\texttt{r\_embed}) is exposed through
        \texttt{trainable\_modules()}, ensuring the optimiser, device transfer,
        \texttt{eval()} mode, and checkpoint coverage all include it.
\end{itemize}

\section{Evaluation Infrastructure}

FID and IS are computed via TorchMetrics. The implementation:
\begin{itemize}
  \item Batches 1{,}000-image evaluation generation in groups of 32 to prevent a single
        large GPU allocation.
  \item Moves completed batches to CPU immediately, bounding VRAM usage to one
        generation batch at a time.
  \item Reuses the real-image reference statistics across all NFE values in one
        evaluation sweep, rather than recomputing per NFE.
  \item Saves cache metadata (sample count, resolution, channels) alongside the
        \texttt{.npz} statistics; a mismatch raises explicitly.
\end{itemize}

\section{UI and Inference Server}\label{sec:ui}

\texttt{web/inference\_server.py} auto-discovers all run directories under
\texttt{results/}, loading checkpoint paths and metric logs at start-up. It
serialises GPU inference requests with a \texttt{threading.Lock} so concurrent
browser sessions do not corrupt GPU state. The training dashboard maps a slider
to checkpoint epochs, loads the corresponding sample image, and overlays the
loss curve truncated at that epoch. Epoch 0 is conceptual noise; epoch 10 is the
first saved checkpoint image.

\section{Reproducibility}\label{sec:reproducibility}

Every run directory contains: a config snapshot, per-epoch JSONL metrics, sample
grids for each NFE, FID reference cache with metadata, and checkpoint archives.
Earlier README text lagged behind the implementation. Before publication it was
updated with fresh-clone setup, Windows/Linux commands, dependency ordering, and
checkpoint-archive guidance. The reproducibility checklist is provided in
Appendix~\ref{app:reproducibility}.
